




Let $E$ and $F$ be fields. 
Recall that $E$ is an extension of $F$ if $F$ is a subfield of $E$.
In this case, $E$ is a vector space over $F$, which we denote as $E/F.$
The degree of $E/F$ is the dimension of $E/F$.
An example of a field extension is a field of the form $E = F(\alpha)$, 
where $\alpha$ is any element over $F.$
Here, $F(\alpha)$ is the smallest field which contains $F$ and $\alpha.$
If $\gamma \in E$ satisfies a polynomial equation with coefficients in $F,$ 
then $\gamma$ is called algebraic over $F$.
The minimum polynomial of an element $\gamma$ in $E$ over $F$ is the monic  
polynomial in $F[X]$ of lowest degree which has $\gamma$ as a root. 
A field extension $E/F$ is algebraic 
if every element in $E$ is algebraic over $F$. 
In particular, $F(\alpha)$ is algebraic over $F$ if $\alpha$ is. 
The degree of $F(\alpha)/F$ is the degree 
of the minimum polynomial of $\alpha$ over $F$.

\bigskip


In this section, we will consider algebraic extensions of finite fields. 
If $F = \bbF_q$ is a field of order $q$, 
then any algebraic extension $E$ 
of $F$ has order $q^e$ where $e$ is the degree of $E$ over $F$.
If $E = F(\alpha)$ is algebraic, 
the degree of $E$ over $F$ 
is the degree of the minimum polynomial of $E$ over $F$.
If $F = \bbF_q$ and $E = F(\alpha)$ 
is algebraic of degree $e$, then $|F| = q^e.$
Every finite field $E$ is of this form, 
where $F = \bbF_p$ and $p$ is the characteristic of $E$.
So, any algebraic extension $E$ of 
$F$ can be constructed as a polynomial factorring 
of the ring $\bbF_p[X].$
For a polynomial $m(X)$ we consider the ideal 
$$
\bI(m) = m(X)\bbF_p[X] = \{m(X)k(X) \mid k(X) \in \bbF_p[X] \} 
$$
of all polynomial multiples of $m(X).$
Under the assumption that $m(X)$ 
has degree $e > 1$ and is irreducible, 
the residue class ring 
$$
\bbF_p[X]/ \bI(m)
$$
is a field with $q = p^e$ elements.
Each residue class has a canonical representative. 
The canonical representative is the unique element 
in the residue class modulo $\bI(m)$
which has degree less than $e$ and leading coefficient one. 

\bigskip


So, for instance, for $q = 4 = 2^2,$ 
we can pick the irreducible polynomial 
$m(X) = X^2+X+1$ over $\bbF_2.$ 
We have four standard representatives modulo $\bI(m),$ namely
$$
\begin{array}{c}
0, \\
1, \\
X, \\
X + 1.\\
\end{array}
$$
Together, these make up a complete set 
of representatives of the residue classes modulo  $\bI(m),$ 
and hence can be identified with the elements of $\bbF_4$:
$$
\bbF_4 = \{0, \; 1, \; X, \; X+1 \}.
$$
The addition of polynomials is as in $\bbF_2[X],$ so 
$$
\begin{array}{c|cccc}
 & 0 & 1 & X & X + 1 \\
\hline
0 & 0 & 1 & X & X + 1\\
1 & 1 & 0 & X+1 & X \\
X & X & X + 1 & 0 & 1 \\
X + 1 & X + 1 & X & 1 & 0 \\
\end{array}
$$
To compute the multiplication table for the field $\bbF_4$.
We can use polynomial arithmetic modulo $m(X):$
It is clear how multiplication 
by $0$ or $1$ works, so we need to focus on the 
polynomials $X$ and $X+1$:
$$
\begin{array}{lllll}
X &\cdot\; X &= X^2 &\equiv X + 1 &\mod X^2 + X + 1, \\
X &\cdot\; (X+1) &= X^2+X &\equiv 1 &\mod X^2 + X + 1, \\
(X+1) &\cdot\; X &= X^2+X &\equiv 1 &\mod X^2 + X + 1, \\
(X+1) &\cdot\; (X+1) &= X^2+1 &\equiv X &\mod X^2 + X + 1, \\
\end{array}
$$
The multiplication table of $\bbF_4$ is
$$
\begin{array}{c|cccc}
 & 0 & 1 & X & X + 1 \\
\hline
0     & 0 & 0   & 0 & 0\\
1     & 0 & 1   & X & X+1 \\
X     & 0 & X   & X+1 & 1 \\
X + 1 & 0 & X+1 & 1 & X \\
\end{array}
$$

\bigskip

Table~\ref{tab:finitefield:activities:more} 
lists Orbiter activities for finite fields. 
This extends Table~\ref{tab:finitefield:activities} in Section~\ref{sec:extension:fields}.  
\orbitertableofcommands{finite_field_activities_2}


\bigskip


The isomorphism type of the resulting field 
only depends on the order $q$ of the field, 
and not on the choice of the polynomial. 
However, for practical computations, the choice of the polynomial matters. 
For instance, results can only be shared between 
different computer algebra systems if the same polynomials are used.
Orbiter has a large collection of polynomials built in. 
Besides these, a polynomial can be specified.
The polynomials that Orbiter offers are in fact primitive, which means that the 
root $\alpha$ is a primitive element for the field $\bbF_q.$ 
This just means that it is a generator of the multiplicative group. 
So, any non-zero element in $\bbF_q$ 
is a suitable power of $\alpha.$


\bigskip

If $\bbF_q$ is an extension of the prime field $\bbF_p$, 
we use a different labeling.
This time, we exploit the fact that $\bbF_q$ is a vector space over $\bbF_p.$
Let $\alpha$ be a root of the irreducible polynomial 
$m(X)\in \bbF_p[X]$ used to create the field. 
Suppose that $q=p^e,$ i.e., the degree of $m(X)$  is $e$.
An $\bbF_p$-basis for the vector space $\bbF_q$ over $\bbF_p$ is given by the powers 
$\alpha^i,$ for $0 \le i < e.$ 
Therefore, any element $\gamma$ of $\bbF_q$ 
has a unique expression of the form 
$$
\gamma = \sum_{h=0}^{e-1} a_i \alpha^i, \quad  0 \le a_i < p  \;\mbox{for all} \; i.
$$
The associated integer rank of $\gamma$ 
is obtained by replacing $\alpha$ by $p$ in this expression 
and evaluating the expression over the integers. 
So, the rank of $\gamma$ is 
$$
\sum_{h=0}^{e-1} a_i p^i.
$$
As $\gamma$ ranges over all field element in $\bbF_q,$ 
the rank values take on every value in the interval $[0,q-1].$ 
The ordering of elements of $\bbF_q$ according 
to these ranks is called the lexicographical ordering. 
The numerical rank of zero is $0$ and the numerical rank of one is $1$.
The numerical rank of $\alpha,$ the primitive element, is $p.$
The numerical ranks of the elements of 
the prime subfield are exactly the elements of $[0,p-1].$

\bigskip

The primitive polynomials used by Orbiter to create small finite fields 
are listed in Table~\ref{tab:primitive:polynomials}. 
The relation is given using the Greek letter that 
is used in orbiter cheat sheets for that particular field.

\orbitertableofextracommands{table_of_irreducible_polynomials}

\bigskip

Let us look at a few examples. 
The command 

\sourcecode{F_4}

creates a report for the field $\bbF_4.$

\bigskip


The command  

\sourcecode{F_4_tables}

creates the addition and multiplication tables for the field:
\orbiteroutput{GRAPHICS/WRAPPERS/F4_tables_add_mult}


\bigskip


The command 

\sourcecode{F_16}

creates a cheat sheet for $\bbF_{16}.$
\orbiteroutput{GRAPHICS/WRAPPERS/cheat_sheet_GF16}


\bigskip

The command  

\sourcecode{F_16_tables}

creates the addition and multiplication tables for the field:
\orbiteroutput{GRAPHICS/WRAPPERS/cheat_sheet_GF16_add_mul_draw}


\bigskip





Unlike  GAP~\cite{GAP} and Magma~\cite{Magma},
Orbiter does not use 
Conway polynomials to create field extensions. 
Instead, it  
provides the option to override the polynomial used to create the finite field.
For subfield relationships, the cheat sheet will indicate the irreducible polynomials 
of all subfields for a given field.  
The table below is taken from a cheat sheet for the field $\bbF_{64}$ 
generated by the polynomial $X^{6} + X^{5} + 1$,  whose numerical rank is 97:
\orbiteroutput{GRAPHICS/LATEX/subfields_of_GF64}

\bigskip

The lexicographic ordering has an interesting 
side-effect for the ordering of elements in extension fields. 
The elements of the prime subfield are always listed before 
any other elements in the extension field. 
For this reason, the addition and multiplication tables of the extension field contain the 
respective table of the prime field in the upper left corner, 
as shown below for the field $\bbF_9$:
\orbiteroutput{GRAPHICS/WRAPPERS/GF_q9_tables_draw}
Here, we omit the zero element in the multiplication table.


\bigskip

Orbiter uses primitive polynomials for creating extension fields. 
Because of this, 
the element $\alpha$ is always primitive. 
Since the numerical rank of $\alpha$ is $p,$ this means that the rank $p$ 
always represents a primitive element in an extension field.
For the addition and multiplication tables of $\bbF_{9}$ 
arranged with respect to powers of a primitive element, see below:
\orbiteroutput{GRAPHICS/WRAPPERS/GF_q9_tables_reordered_draw}

\bigskip


The command  

\sourcecode{nth_roots_63}

creates the minimum polynomials of all elements in $\bbF_{64}^\times.$ 
This corresponds to all irreducible polynomials over $\bbF_2$ 
of degree dividing $6,$ apart from the polynomial $X.$
\orbiteroutput{GRAPHICS/LATEX/Nth_roots_q2_n63}

\bigskip

The command  

\sourcecode{F8_F2_field_reduction}

performs field reduction on the elements $0,\ldots,7$ in $\bbF_8$ down to $\bbF_2.$ 
Here is a graphical representation of the matrices arising this way:
\orbiteroutput{GRAPHICS/WRAPPERS/field_reduction_F8_to_F2}



\bigskip

Let us discuss field embeddings in Orbiter. 
This is important for working in field extensions, 
as there is a restriction on the polynomials used to create the fields.
Suppose we have fields $\bbF_q$ and $\bbF_Q$
with
$$
\bbF_q \le \bbF_Q,
$$
i.e. 
$$
Q = q^e
$$
for some positive integer $e.$
Suppose $\alpha$ is a primitive element 
in $\bbF_Q$ and $\beta$ is a primitive element in $\bbF_q.$
Then Orbiter assumes that  
$$
\beta = \alpha^{\frac{Q-1}{q-1}},
$$
where 
$$
\frac{Q-1}{q-1}
$$
is the index of $\bbF_q^\times$ 
in the multiplicative group $\bbF_Q^\times.$ 
With respect to this embedding, 
the minimum polynomial of $\beta$ is determined. 
Let us consider and example.
The command  

\sourcecode{F_4096}

creates the field $\bbF_{2^{12}}$. 
Orbiter chooses the polynomial 
$$
X^{12} + X^{6} + X^{4} + X + 1
$$
to create the field.
This field has non-trivial subfields $\bbF_4$, $\bbF_8$, $\bbF_{16}$ and $\bbF_{64}.$
The report shows the polynomials for the primitive elements of the subfields:
\orbiteroutput{GRAPHICS/WRAPPERS/table_subfields_GF64}
For each subfield, including the prime field, 
Orbiter computes a subfield basis and the embedding of the elements in the subfield. 
We use the convention that $d=\frac{Q-1}{q-1}.$
Here is the output from the report:
\orbiteroutput{GRAPHICS/WRAPPERS/subfields_GF64_detailed}




\bigskip




The command  

\sourcecode{F_4489}

creates the field $\bbF_{67^2}$. For the sake of speed, 
the option \verb'-without_tables' is given to prevent the field tables from being computed. 


\bigskip

